\section{Finite tensor constructions}
\label{alg:section}

The elementary construction in this section already gives
\(\omega\leq 9\log 3/\log 75=2.290107196\ldots\).
Its main ingredient separates a label that two parties can read independently.
The operation supplies the missing label to the third party and leaves a
large dot product.  We pay for the auxiliary tensor that performs this
operation.  At the exponential scale, the new direct labels compensate for
that payment, while the dot product remains available for matrix
multiplication.  We first prove the operation, including its action on
unresolved tensor factors, and then give the finite constructions used later
for the boundary estimates.

\subsection{Pairing conventions and the input rank}
\label{alg:conventions}

We use the tensor and matrix-dimension conventions of
Section~\ref{sec:preliminaries}.
A pairing on \(X,Y\) of length \(m\), with a scalar third side, is
\(\langle1,m,1\rangle\); the other orientations are
\(\langle m,1,1\rangle\) on \(X,Z\) and
\(\langle1,1,m\rangle\) on \(Y,Z\).
Their tensor products satisfy
\begin{equation}
 \langle a,1,1\rangle\otimes\langle1,b,1\rangle
 \otimes\langle1,1,c\rangle\cong\langle a,b,c\rangle.
 \label{alg:pairing-product}
\end{equation}
The logarithmic volumes of such factors therefore add.

The input is the six-term specialization of the Coppersmith--Winograd
tensor \citep[Section~7, equation~(10)]{cw},
\begin{equation}
 T=CW_1=x_2y_0z_0+x_0y_2z_0+x_0y_0z_2
        +x_0y_1z_1+x_1y_0z_1+x_1y_1z_0.
 \label{alg:cw}
\end{equation}
Write
\(\mathcal S=\{200,020,002,011,101,110\}\) for its support, abbreviating
\((i,j,k)\) by \(ijk\).  At every position the three indices sum to two.
The maps in our square construction for this input have algebraic
coefficients.  The general preservation identities below also allow
arbitrary complex coefficients in the unresolved tensors.

\begin{lemma}
\label{alg:cw-rank}
The tensor \(T\) has border rank three.  Consequently
\(\brank(T^{\otimes n})\leq3^n\) for every integer \(n\geq1\).
\end{lemma}
\begin{proof}
Let \(\zeta=e^{2\pi i/3}\), and put
\[
 X_r(t)=x_0+\zeta^r t x_1+\zeta^{2r}t^2x_2
 \quad (r=0,1,2),
\]
with analogous definitions for \(Y_r,Z_r\).  Filtering the degree modulo
three gives the exact identity
\[
 \frac{1}{3t^2}\sum_{r=0}^2\zeta^r X_r(t)Y_r(t)Z_r(t)
 =T+t^3(x_2y_2z_1+x_2y_1z_2+x_1y_2z_2).
\]
For every nonzero \(t\), the left side has rank at most three.  Taking
\(t\to0\) proves the upper bound.  In the flattening with the \(X\) side
as rows, the three coefficients of \(x_0,x_1,x_2\) are linearly independent:
the monomials \(y_0z_2,y_0z_1,y_0z_0\), respectively, distinguish them.
The flattening rank is therefore three, giving the reverse bound.  Tensoring
the upper-bound construction proves the last assertion.
The displayed identity also supplies an explicit polynomial witness of
cost three in the terminology of Section~\ref{sec:preliminaries}.
\end{proof}

We use the following established theorem in exactly its border-rank form;
see also the explicit statement in
\citet[Theorem~3.1]{afl}.

\begin{theorem}[Asymptotic sum inequality, \citet{schonhage}]
\label{thm:asi}
If a tensor \(U\) degenerates to
\(\bigoplus_{\ell=1}^{L}\langle a_\ell,b_\ell,c_\ell\rangle\), with all
dimensions positive integers, then
\[
 \sum_{\ell=1}^{L}(a_\ell b_\ell c_\ell)^{\omega/3}
 \leq\brank(U).
\]
In particular, if
\(T^{\otimes n}\otimes\mathcal A\unrhd
\bigoplus_{\ell=1}^{L}\langle a,b,c\rangle\)
and \(\brank(\mathcal A)\leq R\), then
\begin{equation}
 L(abc)^{\omega/3}\leq3^nR.
 \label{alg:asi-account}
\end{equation}
\end{theorem}

The auxiliary rank \(R\) in \eqref{alg:asi-account} is part of the input
cost.  Thus, for a proposed exponent \(\beta\), a finite realization with
\begin{equation}
 \log L-\log R+\frac{\beta}{3}\log(abc)>n\log3
 \label{alg:strict-score}
\end{equation}
proves \(\omega<\beta\).  We never cancel an auxiliary tensor as a tensor
identity; only its explicitly charged logarithmic rank will be offset by
output multiplicity.

\subsection{Supplying a label to its missing party}
\label{alg:group-section}
\label{sec:separation}

Suppose a tensor is partitioned into pieces whose \(X\)-variables are
disjoint and whose \(Y\)-variables are disjoint.  Each of these two parties
can then determine the same piece label from its own variable.  The
\(Z\)-variables may still be shared.  The next lemma resolves this last
sharing without resolving the terms inside a piece.  We call this operation
GROUP.

The construction uses equal-norm shells, following the geometry of
\citet{behrend}, together with affine slices that constrain two independently
chosen auxiliary indices.  Its auxiliary tensor is the multiplication
tensor of a finite abelian group, linking the construction to the
group-theoretic tensor approach of \citet{cohn-umans}.  The identity and its
full rank charge are proved below.

\begin{lemma}[Pairwise label separation]
\label{lem:group}
\label{lem:separator}
Let \(U=\sum_{s=1}^{K}U_s\), where \(U_s\) uses \(X\)-variables in
\(X_s\) and \(Y\)-variables in \(Y_s\), and the families
\((X_s)_{s=1}^K\) and \((Y_s)_{s=1}^K\) are each disjoint.  The
\(Z\)-variables may overlap.  There are a finite abelian group \(G\), its
group multiplication tensor \(\mathcal A_G\), and a positive integer
\(m\) such that
\begin{equation}
 U\otimes\mathcal A_G\unrhd
 \bigoplus_{s=1}^{K}\bigl(U_s\otimes\langle1,m,1\rangle\bigr).
 \label{alg:group-identity}
\end{equation}
The degeneration is a zeroing-out on auxiliary coordinates controlled by
the original local variables.  It preserves all original coordinates and
coefficients.  For \(K\to\infty\), the parameters can be chosen with
\begin{equation}
 \rank(\mathcal A_G)=\brank(\mathcal A_G)=|G|,
 \qquad
 \log|G|=\log K+O(\sqrt{\log K}),
 \qquad
 m=K.
 \label{alg:group-rates}
\end{equation}
The same statement holds for any other pair of knowing sides.
\end{lemma}
\begin{proof}
The case \(K=1\) uses the trivial group and \(m=1\), so no
nontrivial auxiliary tensor is needed.
For \(K\geq2\), set
\[
 d=\lceil\sqrt{\log K}\rceil+3,
 \qquad
 Q=\left\lceil((d+1)K)^{1/(d-2)}\right\rceil,
 \qquad m=K.
\]
Since \(dQ^2+1\leq(d+1)Q^2\), these choices give
\(Q^d/(dQ^2+1)\geq K\).
The grid \(\{0,\ldots,Q-1\}^d\) has \(Q^d\) points and at most
\(d(Q-1)^2+1\) possible squared norms.  Pigeonholing supplies at least
\(K\) distinct points \(t_1,\ldots,t_K\) on one Euclidean sphere.
Assign \(t_s\) to label \(s\).

For each \(s\), the integer inner product \(t_s\cdot u\), as \(u\)
ranges over this grid, has at most \(d(Q-1)^2+1\) values.  A further
pigeonhole argument supplies a set \(V_s\) of at least \(m\) grid points
on one affine slice
\[
 t_s\cdot u=h_s.
\]
Trim each \(V_s\) to exactly \(m\) points.  Only the common inner product
within a slice is needed; the points of a slice need not have equal norm.

Take
\[
 G=(\Z/(4Q)\Z)^d,
 \qquad
 \mathcal A_G=\sum_{a+b=c}\mathsf{x}_a\mathsf{y}_b\mathsf{z}_c,
\]
using separate auxiliary variables on the three sides.  For completeness,
character orthogonality gives a rank decomposition
\[
 \mathcal A_G=\frac1{|G|}\sum_{\chi\in\widehat G}
 \left(\sum_{a\in G}\chi(a)\mathsf{x}_a\right)
 \left(\sum_{b\in G}\chi(b)\mathsf{y}_b\right)
 \left(\sum_{c\in G}\chi(-c)\mathsf{z}_c\right).
\]
It has \(|G|\) summands.  The \(X\)-flattening has \(|G|\) linearly
independent rows, since the coefficient of
\(\mathsf{y}_0\mathsf{z}_a\) distinguishes row
\(a\).  This proves both rank equalities.

For an original \(X\)-variable in \(X_s\), retain auxiliary coordinates
\(a=u\in V_s\).  For an original \(Y\)-variable in \(Y_s\), retain
\(b=t_s-u'\) with \(u'\in V_s\).  On the \(Z\) side retain only
\(c\in\{t_1,\ldots,t_K\}\).  Negative coordinates of \(t_s-u'\)
are interpreted modulo \(4Q\).  These are three local variable
projections: each knowing side uses the label determined by its original
variable, and the third side uses only the prescribed tag set.

A surviving term satisfies
\(c\equiv t_s+u-u'\pmod{4Q}\).
Each coordinate of \(c-t_s-u+u'\) has absolute value at most
\(2(Q-1)<4Q\), so the congruence is an equality of integer vectors.
The tag norms are equal and
\(t_s\cdot(u-u')=0\), whence
\[
 0=\|c\|^2-\|t_s\|^2=\|u-u'\|^2.
\]
Thus \(u=u'\) and \(c=t_s\).  Conversely, every \(u\in V_s\) gives
one surviving triple \((u,t_s-u,t_s)\).
The auxiliary factor in piece \(s\) is precisely
\[
 \sum_{u\in V_s}\mathsf{x}_u\mathsf{y}_{t_s-u}\mathsf{z}_{t_s},
\]
a length-\(m\) pairing on the knowing sides.  The pieces were already
disjoint on those sides and now have distinct auxiliary tags on the third.
The calculation also excludes every cross term, proving
\eqref{alg:group-identity}.

Finally, with \(L=\log K\), the displayed choice of \(d,Q\) satisfies
\[
 d\log Q
 \leq\frac{d}{d-2}\bigl(L+\log(d+1)\bigr)+o(1)
 =L+O(\sqrt L).
\]
The rounding error in \(\log Q\) is exponentially small in
\(\sqrt L\).  Since \(K\leq Q^d\) and
\(|G|=(4Q)^d\), these inequalities and \(m=K\) prove
\eqref{alg:group-rates}.
\end{proof}

For repeated uses at different label counts, fix one such group \(G_K\)
for each \(K\), and write \(A_K=\mathcal A_{G_K}\).  This is the
auxiliary-tensor notation used in the conditional-label entropy bound.

The geometry rules out a stronger collision than the usual
three-term-progression condition: the two slice points \(u,u'\) are chosen
independently.  Their difference is perpendicular to \(t_s\); any nonzero
such displacement from \(t_s\) lies outside its sphere.  Requiring the
endpoint to be another tag therefore forces that difference to vanish.
Figure~\ref{fig:group-geometry} shows this displacement argument.

\begin{figure}[H]
\centering
\begin{tikzpicture}[>=Latex,scale=1.0,font=\small]
\draw[blue!60!black,thick] (0,0) circle (1.9);
\draw[gray!55,dashed] (1.9,-2.1)--(1.9,2.4);
\fill (0,0) circle(1.5pt) node[below left] {\(0\)};
\coordinate (tag) at (1.9,0);
\coordinate (candidate) at (1.9,1.65);
\draw[->,thick] (0,0)--(tag) node[midway,below] {\(t_s\)};
\draw[->,thick,orange!75!black] (tag)--(candidate) node[midway,right] {\(d=u-u'\)};
\draw[->,thick,gray!75!black] (0,0)--(candidate) node[midway,above left] {\(c=t_s+d\)};
\fill (tag) circle(2pt) node[below right] {\(t_s\)};
\fill (candidate) circle(2pt) node[above right] {\(c\)};
\draw (1.65,0)--(1.65,.25)--(1.9,.25);
\node[align=left,anchor=west,text width=6cm] at (4.1,0.6) {
The affine slice gives \(t_s\cdot d=0\).\\[5pt]
The group equation gives \(c=t_s+d\).\\[5pt]
If \(c\) is another tag on the same sphere,\\[3pt]
\(\|c\|^2=\|t_s\|^2+\|d\|^2=\|t_s\|^2\),\\[3pt]
so \(d=0\) and \(c=t_s\).
};
\end{tikzpicture}
\caption{The separation identity in a two-dimensional slice.
The dashed line represents the displacement hyperplane
\(t_s+t_s^\perp\), rather than the slice containing \(u,u'\).
A nonzero displacement along it leaves the tag sphere, so no other tag
can survive. The argument applies in every dimension.}
\label{fig:group-geometry}
\end{figure}


\begin{lemma}[Preservation of unresolved fibers]
\label{alg:payload}
The identity in Lemma~\ref{lem:group} holds with arbitrary coefficients
and arbitrary unresolved local indices inside the tensors \(U_s\),
including indices correlated between different original occurrences.
Local support projections and monomial leading-term restrictions have the
same property.  Finite compositions of these operations may leave all
previously produced matrix multiplication factors untouched.
\end{lemma}
\begin{proof}
Retain an original basis coordinate on every side throughout the
construction.  In Lemma~\ref{lem:group}, a surviving auxiliary triple is
characterized solely by the piece label and the equation just proved.
The original coefficient is neither changed nor combined with another
coefficient.  Appending additional local indices to any original variable
therefore gives the same identity for every coefficient separately.  The
extra indices may be arbitrary; no product distribution or independence
between occurrences is used.

A local support projection has the same coefficientwise description.
For a monomial restriction, the degree of each term is the sum of three
local degrees, and its leading coefficient is either its original
coefficient or zero.  This too commutes with adjoining unresolved
coordinates or earlier local count predicates.  The claim for a finite
composition follows by induction.  At each subsequent stage, tensor the
new maps with the identity on all previously produced pairings.
\end{proof}

This preservation statement is a locality statement, not a nonvanishing
statement.  If an earlier support restriction makes a designated source
word absent, its fiber is zero and cannot be counted as an output.  Our
finite constructions below explicitly retain every counted word.  The
later count construction will require a separate guarantee for every
complete controlled path.

\subsection{Exact types and a finite hierarchy of labels}
\label{alg:hierarchy-section}

For a probability vector \(p\) on a finite alphabet, write
\(H(p)=-\sum_i p_i\log p_i\), with \(0\log0=0\).
Exact types provide the finite counting form of entropy used here;
see \citet{csiszar1998} for the method of types.
For rational \(p\) and admissible integers \(N\), the number of words
having exactly \(Np_i\) occurrences of symbol \(i\) is
\begin{equation}
 \binom{N}{(Np_i)_i}
 =\exp\bigl(NH(p)+O(\log(N+1))\bigr).
 \label{alg:type-count}
\end{equation}
The implicit constant depends only on the fixed alphabet.  We use exact
types so that the number of positions assigned to every later conditional
operation is independent of the previously selected direct output.

The following lemma makes precise how multiple separation steps share
auxiliary tensors and count each label once.  Its accounting is the
conditional-entropy chain rule \citep{shannon1948}, applied to labels
that the tensor maps can actually read.

\begin{lemma}[Accounting for a finite label hierarchy]
\label{lem:hierarchy}
Let \(b\geq1\), and let \(\mathcal B\subseteq\mathcal S^b\) be
nonempty.  Suppose local restrictions and degenerations of
\(T^{\otimes b}\) retain exactly these original scalar components, with
their original coordinates and coefficients.  Suppose there are finitely many
labels \(S_1,\ldots,S_t\), each a function of a component in
\(\mathcal B\), with these properties:
\begin{enumerate}[label=\textup{(\roman*)}]
 \item the complete record \((S_1,\ldots,S_t)\) determines the original
       component;
 \item after the earlier labels are fixed, the next label is independently
       readable by two fixed sides within that conditional pool;
 \item every component designated by a retained complete record survives
       all the local restrictions.
\end{enumerate}
Let \(\alpha\) be a rational probability distribution on \(\mathcal B\).
For admissible \(N\to\infty\), there are auxiliary tensors
\(\mathcal A_N\) and degenerations of
\(T^{\otimes bN}\otimes\mathcal A_N\) into \(L_N\) independent matrix
multiplication tensors of common volume \(v_N\), with
\(\brank(\mathcal A_N)\leq R_N\), such that
\begin{equation}
 \log L_N=NH(\alpha)+o(N),\qquad
 \log R_N=NH(\alpha)+o(N),\qquad
 \log v_N=NH(\alpha)+o(N).
 \label{alg:hierarchy-rates}
\end{equation}
When \(H(\alpha)>0\), this gives
\(\omega\leq3b\log3/H(\alpha)\).
\end{lemma}
\begin{proof}
Apply the stipulated local restrictions in each of the \(N\) blocks.
At layer \(j\), write \(h=(S_1,\ldots,S_{j-1})\) for a preceding
record, let \(p_h\) be its probability under \(\alpha\), and let
\(p_{s\mid h}\) be the conditional law of \(S_j\).
The earlier layers have made the full preceding-label word direct on all
three sides.  In every branch there are exactly \(Np_h\) block positions
in conditional pool \(h\).  Each knowing side can read the next-label
word in those positions.  It can therefore zero variables whose next-label
word does not have counts \(Np_hp_{s\mid h}\).  This is a genuine local
type projection, not a presumed projection onto an unavailable joint type.

There remain
\[
 K_{j,h}=\binom{Np_h}{(Np_hp_{s\mid h})_s}
\]
possible next-label words in this pool.  Use Lemma~\ref{lem:group} to
separate them.  A zero-probability pool is omitted, and a deterministic
next label needs no auxiliary tensor.  For a positive-entropy pool,
\eqref{alg:type-count} and \eqref{alg:group-rates} give identical leading
logarithmic contributions
\[
 Np_h H(S_j\mid h)+o(N)
\]
to its direct count, auxiliary rank, and pairing volume.

The auxiliary tensor for pool \((j,h)\) is appended once.  It is reused
in every preceding direct branch through a block-diagonal map: on a direct
sum, each of the three parties knows the branch, so it may address the
appropriate positions using that branch's label.  Equivalently,
\((\bigoplus_\ell V_\ell)\otimes\mathcal A\) is the direct sum of the
\(V_\ell\otimes\mathcal A\), and separate maps may act in its blocks.
This does not charge \(\rank(\mathcal A)\) once per \(\ell\).
The size and physical orientation of every conditional supply depend only
on \((j,h)\), whose population is fixed, not on the selected output word.

There are only finitely many layers and conditional pools.  Multiplying
their counts, ranks, and pairing volumes, and using the entropy chain rule,
gives in each case the leading logarithm
\[
 N\sum_{j=1}^t H(S_j\mid S_1,\ldots,S_{j-1})=NH(\alpha).
\]
All error terms sum to \(o(N)\).  The complete record determines one
original scalar component in each block, and every such designated
component is nonzero by assumption.  Exact conditional counts fix the
number of pairings in each physical orientation.  Formula
\eqref{alg:pairing-product} therefore gives a common matrix multiplication
tensor in all final branches.  This proves \eqref{alg:hierarchy-rates}.
Substituting them in \eqref{alg:asi-account}, dividing by \(N\), and taking
the limit proves the exponent bound.
\end{proof}

Nothing in this lemma permits arbitrary support deletion.  The locally
obtainable support and the two-sided readability must be established for
each hierarchy.  Once they are established, exact conditional types select
its desired joint frequencies through accessible labels.

\subsection{An explicit four-component construction}
\label{alg:four}

First remove the term \(101\) from \(T\).  Scale each of
\(x_1,z_1,y_0,y_2\) by \(t\).  The five other terms have degree one in
\(t\), whereas \(101\) has degree three.  Dividing by \(t\) and taking
the limit gives
\begin{equation}
 T\unrhd T'=A+B+C,
 \quad
 A=x_0(y_0z_2+y_1z_1),\quad
 B=(x_2y_0+x_1y_1)z_0,\quad
 C=x_0y_2z_0.
 \label{alg:five}
\end{equation}
A local predicate on each side identifies exactly one group:
\(Z>0\) identifies \(A\), \(X>0\) identifies \(B\), and \(Y=2\)
identifies \(C\).  Exactly one of these predicates holds on every retained
term.  We call any such three-group partition a \emph{W-partition}; the
three predicates, rather than the original numerical indices, are its
one-hot flags.

Delete \(C\) by setting \(y_2=0\), and take \((A+B)^{\otimes N}\),
with \(4\mid N\).  Retain coarse words with \(N/2\) copies of \(A\)
and \(N/2\) copies of \(B\).  This is local on \(X\), since a zero
\(X\) index means \(A\) and a positive one means \(B\).
The \(Z\) side also identifies the word, with the roles reversed.
There are \(\binom{N}{N/2}\) coarse words.  Lemma~\ref{lem:group}
separates them by supplying their label to \(Y\), and gives each word a
pairing on \(X,Z\) of logarithmic length \(N\log2+o(N)\).

Inside a fixed coarse word, \(A\)'s two terms \(002,011\) are each
identified by \(Y\) and by \(Z\).  Retain \(N/4\) occurrences of each
among the \(N/2\) \(A\) positions and separate their words.  This adds a
pairing on \(Y,Z\) of logarithmic length \((N/2)\log2+o(N)\).
Similarly, \(B\)'s two terms \(200,110\) are each identified by
\(X,Y\); resolving their equal-frequency words adds a pairing on
\(X,Y\) of the same logarithmic length.  These restrictions are
conditional on the already-direct coarse word and therefore are local.

The three pairings combine by \eqref{alg:pairing-product}.  The total
logarithmic output count, auxiliary rank, and volume are each
\(N\log4+o(N)\).  Theorem~\ref{thm:asi} consequently gives
\begin{equation}
 \omega\leq\frac{3\log3}{\log4}=2.377443751\ldots.
 \label{alg:four-bound}
\end{equation}
The scalar original terms alone would not supply any volume to which
\(\omega\) could be applied.  Here the three pairing orientations supply
that volume, while all three auxiliary rank costs have been included.

\subsection{Completion and the 75-component bound}
\label{alg:completion-section}

A W-partition has a useful local description even when its three groups
contain unresolved tensors: side \(i\) knows precisely the positions at
which its own group occurs.  This permits a degeneration of a block of
several occurrences before any group has been made direct.  We call the
following operation COMPLETE.

\begin{lemma}[Completion]
\label{alg:completion}
Let a tensor have a W-partition with labels \(0,1,2\).  For an integer
\(b\geq1\), a local monomial degeneration of its \(b\)-th power retains
exactly the coarse words
\begin{equation}
 \{0,1\}^b\ \cup\ \{0,2\}^b.
 \label{alg:complete-support}
\end{equation}
This support has a new W-partition: the constant word \(0^b\), the words
containing \(1\), and the words containing \(2\).  Conditional on either
nonzero new label, the original binary word is independently readable by
two sides.  The construction preserves arbitrary unresolved fibers.
\end{lemma}
\begin{proof}
The three local predicates are
\[
 \text{every position has label }0,\qquad
 \text{some position has label }1,\qquad
 \text{some position has label }2.
\]
Each is readable on the correspondingly active side of the original
W-partition.  Their sum is one on \eqref{alg:complete-support} and two
when both \(1\) and \(2\) occur.  Multiply a variable by \(t\) when its
predicate holds, divide the tensor by \(t\), and take the limit.  The
surviving predicates form a new one-hot partition.  In the nonzero group
\(i\in\{1,2\}\), the word uses only \(0,i\), so either of those two
active sides determines the entire word from the positions of its own
label.  The degeneration is coefficientwise and hence preserves the fibers
by Lemma~\ref{alg:payload}.
\end{proof}

For the groups \(A,B,C\) in \eqref{alg:five}, assign labels \(0,1,2\)
in that order.  The group sizes are \(2,2,1\).  A completed block has
three group sizes
\[
 2^b,\qquad4^b-2^b,\qquad3^b-2^b.
\]
Delete its all-\(A\) group: its active \(Z\) predicate is that every
original index is positive.  The remaining groups will be denoted
\(P_B\) and \(P_C\).  The \(X\) side recognizes \(P_B\) by the
occurrence of a positive index, and the \(Y\) side recognizes \(P_C\)
by the occurrence of index two.  Since only these two groups remain,
either side independently determines which group occurs.

\begin{theorem}[Elementary completion bound]
\label{thm:elementary}
For every integer \(b\geq2\),
\begin{equation}
 \omega\leq
 \frac{3b\log3}{\log(4^b+3^b-2^{b+1})}.
 \label{alg:complete-bound}
\end{equation}
In particular,
\[
 \omega\leq\frac{9\log3}{\log75}=2.290107196\ldots.
\]
\end{theorem}
\begin{proof}
Apply \eqref{alg:five} to each occurrence, perform the completion of
Lemma~\ref{alg:completion} in blocks of size \(b\), and delete the
all-\(A\) group as just described.  This locally retains
\(s_b=4^b+3^b-2^{b+1}\) original scalar components per block.
Here is a complete hierarchy that separates them.

First resolve the binary group label \(P_B,P_C\), known on \(X,Y\).
After it is direct, resolve the original word inside each group: the
\(A,B\) word in \(P_B\) is known on \(Z,X\), and the \(A,C\) word
in \(P_C\) is known on \(Z,Y\).  After that word is direct, resolve
\(002\) versus \(011\) in every \(A\) position using \(Y,Z\), and
\(200\) versus \(110\) in every \(B\) position using \(X,Y\).
There is no internal choice in \(C\).  The complete record identifies
one original scalar component.  Every such component survives the stated
monomial restriction and deletion.  The pair of knowing sides at each
stage is fixed by the preceding direct labels.

Give the \(s_b\) components equal frequency and apply
Lemma~\ref{lem:hierarchy}.  Their entropy is \(\log s_b\), while the
original rank cost per block is at most \(3^b\).  This proves
\eqref{alg:complete-bound}.

For \(b=3\), the sizes of \(P_B,P_C\) are \(56,19\).  The prescribed
outer frequencies are \(56/75,19/75\); the internal component law in
each group is uniform.  The entropy identity
\[
 H(56/75,19/75)+\frac{56}{75}\log56+\frac{19}{75}\log19=\log75
\]
shows explicitly how the successive choices count each of the 75 terms
once.  Over \(N\) such blocks the three leading logarithms are
\(\log L_N=\log R_N=\log v_N=N\log75+o(N)\).
Using the input cost \(3^{3N}R_N\) in
\eqref{alg:asi-account} and canceling the \emph{logarithmic} multiplicity
and auxiliary costs leaves
\((\omega/3)\log75\leq3\log3\), as claimed.
\end{proof}

The neighboring block lengths illustrate that this particular completion
is not monotone in \(b\):
\begin{center}
\begin{tabular}{@{}ccc@{}}
\toprule
Block length & Retained scalar components & Exponent bound \\
\midrule
2 & 17 & \(2.326571610\ldots\) \\
3 & 75 & \(2.290107196\ldots\) \\
4 & 305 & \(2.304655406\ldots\) \\
\bottomrule
\end{tabular}
\end{center}
No optimization claim over all restrictions of the corresponding power is
needed.  The length-three construction already gives the stated bound.

\subsection{Parity and conditional W constructions}
\label{alg:parity-section}

Completion gives one way of making a new W-partition.  A different local
restriction on three occurrences supplies two useful W labels, which can
be resolved in either order.  We call this operation PARITY.

\begin{lemma}[Parity restriction]
\label{alg:parity}
In three W-partitioned occurrences, a local monomial degeneration removes
exactly the six words in which all three labels occur and retains the
other 21 coarse words.  On the retained support, the majority label and
the odd-count label each define a W-partition.  Once both labels are direct,
the remaining choice is either deterministic or independently readable by
two sides.  All statements remain true with arbitrary unresolved fibers.
\end{lemma}
\begin{proof}
Let \(n_i\) be the count of label \(i\) among the three positions.
Side \(i\) knows its occurrence set and hence \(n_i\).
Scale a variable by \(t^{n_i\bmod2}\) on that side.
The sum of the three parity exponents is one for count patterns
\((3,0,0)\) and \((2,1,0)\), up to permutation, and three for
\((1,1,1)\).  Division by \(t\) and passage to the limit gives the
claimed support.

On this support exactly one label has count at least two, and exactly one
has odd count.  Both are local unary tests, giving the two W-partitions.
Write these labels as \(c,h\).  If \(c=h\), the word is \(ccc\).
Otherwise, it consists of two \(c\)'s and one \(h\).  The \(c\) side
knows the two majority positions and the \(h\) side knows the singleton
position.  Either determines which of the three placements occurs, so
Lemma~\ref{lem:group} applies after \(c,h\) have been made direct.
All restrictions and the final separation preserve fibers by
Lemma~\ref{alg:payload}.
\end{proof}

If the three occurrences have internal scalar counts
\((a_i)_i,(b_i)_i,(c_i)_i\), the fiber with majority \(i\) and odd-count
label \(j\) has size
\begin{equation}
 M_{ij}=
 \begin{cases}
  a_ib_ic_i,&i=j,\\
  a_ib_ic_j+a_ib_jc_i+a_jb_ic_i,&i\ne j.
 \end{cases}
 \label{alg:parity-weights}
\end{equation}
Each term for \(i\ne j\) describes one singleton placement.
For the boundary weights \((2,2,1)\) on all three occurrences, this is
\[
 (M_{ij})=\begin{pmatrix}8&24&12\\24&8&12\\6&6&1\end{pmatrix}.
\]
For example, delete the third majority row and resolve the remaining
binary majority label.  Within the first direct row delete its first
odd-count column; within the second direct row delete its second column.
These are permitted conditional local deletions because the majority is
already direct.  Each row then has size 36, giving 72 original scalar
components.  Lemma~\ref{lem:hierarchy} yields
\(\omega\leq9\log3/\log72=2.311966920\ldots\).
This illustrative restriction is weaker than Theorem~\ref{thm:elementary}.
Parity is useful because its two conditional groupings enlarge the finite
construction family at nonuniform distributions.

\subsection{Finite recipes and weighted laws}
\label{alg:recipes}

We now track the source law and score of a finite tree of W constructions.
Fix labels \(0,1,2\) on each original W-partitioned occurrence before
applying the tree.  These are its \emph{canonical labels}: later physical
permutations may change which side reads a flag, but the original labels
are still counted in this fixed ordering.

For the score calculations below, fix a finite operation tree and its
rational conditional laws.  Admissible exact-type populations, together
with the efficient GROUP tensors of \eqref{alg:group-rates}, give a family
of finite realizations.  The \emph{limiting score} of this fixed tree and
law data is the limit along increasing populations; each member remains
one finite recipe with fixed tensor supplies and maps.
A later operation may depend on a label only once that label is direct on
all three sides.  At every stage, previously produced pairings are separate
tensor factors; their internal indices are never used to choose subsequent
controls.

For a realization on \(n\) original occurrences, let \(L\) be its direct
multiplicity, \(R\) its auxiliary rank bound, and \(v\) its additional
matrix multiplication volume.  The rates are
\[
 D=\frac{\log L}{n},\qquad A=\frac{\log R}{n},\qquad
 M=\frac{\log v}{n}.
\]
Here \(M\) excludes any tensor initially present inside a coarse group.
For the W recipes the convenient formal score is
\(u_\gamma=\gamma(D-A)+M\), where \(\gamma\geq1\).
For \(0<\beta\leq3\),
\[
 u_{3/\beta}=\frac3\beta\left(D-A+\frac\beta3M\right).
\]
Thus \(u_\gamma\) rescales the normalized score expression in
\eqref{alg:strict-score} so that the coefficient of the added volume is one.
Its definition assumes no bound on \(\omega\).
Definitions~\ref{def:strong} and~\ref{val:def-value} will specify the
banked realizations and the closure of their law--score pairs.
For a fixed tree consisting only of
pairwise separation, completion, parity, and conditional exact types,
its total \(D-A\) tends to zero when the group/type supplies grow.
Consequently its limiting score is independent of \(\gamma\).

A binary leaf deletes one W label locally and resolves the remaining
binary word by Lemma~\ref{lem:group}.  At law \(q\) on that pair, its
limiting score is \(H(q)\).  If positive numbers \(d_i\) are assigned as
weights for choosing a law, the weighted score is
\(H(q)+\sum_iq_i\log d_i\).  Maximizing on the retained pair \(i,j\)
gives \(\log(d_i+d_j)\).  These weights may guide selection of a finite
rational law; they do not change the coefficients of a tensor.

For completion of \(b\) occurrences with weights \(d_{r,i}\), the three
fiber weights are
\begin{equation}
 w_0=\prod_{r=1}^b d_{r,0},\qquad
 w_i=\prod_{r=1}^b(d_{r,0}+d_{r,i})-\prod_{r=1}^b d_{r,0}
 \quad(i=1,2).
 \label{alg:complete-weights}
\end{equation}
When the weights are scalar term counts, this is literal counting of the
supports in Lemma~\ref{alg:completion}.  For arbitrary positive weights
it is the corresponding generating polynomial.  For parity the analogous
matrix is \eqref{alg:parity-weights}.

\subsection{Physical orientations and canonical laws}

A physical permutation acts on the full tensor factors of an occurrence or
completed block.  Each original occurrence keeps both its canonical label
and its net physical side permutation.  For
\(\sigma\in\mathfrak S_3\), the \emph{orientation bank} \(\sigma\)
is the collection of original occurrences with that net permutation
relative to their canonical coordinates.  We will arrange that all six
banks have the same exact canonical label histogram.  Their common
normalized histogram is the source law; it is not the aggregate law of the
physical flags after the bank orientations are forgotten.

The exact types used below fix canonical histograms within each relative
orientation.  The next lemma shows how the six global physical permutations
then distribute these counts equally among the original banks, while
keeping one common auxiliary supply.

\begin{lemma}[Equal physical orientation supplies]
\label{alg:orientation-banks}
Fix a finite oriented schedule of operations with prescribed exact
conditional populations.  Suppose its exact types fix, for each relative
physical permutation \(\tau\in\mathfrak S_3\), the canonical input
histogram \(H_\tau\) of all original occurrences with that permutation,
independently of the direct history.  Include all six global physical side
permutations of this schedule, with the same prescribed populations and
histograms in corresponding replicas.
Then every original orientation bank has the same exact canonical input
histogram.  Auxiliary supplies and pairing dimensions are independent of
the direct history.  All regroupings controlled by preceding direct labels
are implementable by local block-diagonal maps.
\end{lemma}
\begin{proof}
In the replica with global permutation \(\sigma\), an original occurrence
of relative orientation \(\tau\) has absolute orientation
\(\sigma\tau\).  Thus an absolute bank \(\rho\) receives the histogram
\(H_\tau\) from the unique replica
\(\sigma=\rho\tau^{-1}\).  Its total histogram is
\[
 \sum_{\tau\in\mathfrak S_3}H_\tau,
\]
independently of \(\rho\) and of the direct history.  The occurrence
positions contributing to a histogram may vary with the history, but its
counts are fixed by hypothesis.  Each original occurrence is counted once,
in its relative-orientation class.

All parties know a preceding direct label, so each can reorder or address
the same specified positions using its own local coordinate map.
Corresponding relative permutations address corresponding physical banks;
no party is required to read another party's local flag.  Each fixed pool
uses a fixed auxiliary tensor and fixed pairing lengths.  Tensoring all
those supplies once and acting block-diagonally therefore gives uniform
rank bounds and dimensions.  Previously produced pairings remain untouched.
\end{proof}

In particular, rotating the flags of an intermediate completed block
changes its grouping and its readers, but does not recode the labels of
the original occurrences inside it.  To recover a canonical label from a
physically permuted occurrence, undo that occurrence's permutation.

\subsection{Resolving and composing nested recipes}

An operation in the finite library may complete two occurrences, physically
transpose that completed block, and then complete it with a third
occurrence.  Resolving its original word calls for a hierarchy inside the
final fiber.  The following statement supplies that hierarchy for any
fixed finite composition of completion gates.

\begin{lemma}[Resolving a tree of completion gates]
\label{alg:completion-hierarchy}
Consider a finite rooted tree whose leaves are W-partitioned occurrences
and whose internal nodes apply Lemma~\ref{alg:completion} to their children.
Before a child is used at a parent, allow a fixed physical permutation of
that whole child block.  The retained support has a W-partition at its
root.  Once the root label is direct, the original leaf-label word admits
a finite hierarchy of pairwise-readable refinements.  All refinements
preserve arbitrary unresolved leaf fibers.

In a fixed root fiber, let \(P\) be any rational law on the retained
original leaf-label words.  Across \(N\) copies of that fiber, exact
conditional types and Lemma~\ref{lem:group} give direct multiplicity,
auxiliary rank bound, and added pairing volume whose logarithms are each
\(NH(P)+o(N)\).  The auxiliary product and pairing dimensions are common
to every direct output.  Canonical original labels are unchanged by the
intermediate physical permutations.
\end{lemma}
\begin{proof}
Process the nodes from the root toward the leaves.  Suppose a node's label
is already direct, and express its children's W labels in that node's
current physical coordinates.  If the node label is zero, completion
forces all those child labels to be zero.  Otherwise, if its label is
\(i\in\{1,2\}\), its child-label word lies in \(\{0,i\}^b\), where
\(b\) is the number of children.  Side zero knows the positions with label
zero and side \(i\) knows the positions with label \(i\); either side
determines the word.  Apply Lemma~\ref{lem:group} to separate that word.
Every child label is now direct on all three sides.  Undo the known
physical permutation on its label to identify the child's own root label,
and recurse into that child.  Recursion is required even when the parent
label was zero: a known child root may still contain unresolved choices.

There are finitely many nodes.  Each successive refinement is a function
of the original leaf word, and the complete record recovers that word.
For the law \(P\), prescribe the induced rational conditional type of
each refinement in each preceding direct pool.  These types are locally
readable by the two sides just identified.  They fix all pool populations,
so a single prescribed auxiliary product can be used by block-diagonal
maps in every branch.  They also fix the counts of pairings in every
physical orientation.  As in the proof of Lemma~\ref{lem:hierarchy}, each
pool contributes its conditional entropy to all three leading logarithms.
The chain rule sums these contributions to \(H(P)\), since the complete
record and the original leaf word determine one another.  The output
pairing dimensions are therefore common to all branches.

The retained support is exactly the intersection of the primitive
completion conditions, each implemented by its local monomial
degeneration.  Every word in \(P\)'s support satisfies those conditions.
Lemma~\ref{alg:payload} preserves any unresolved leaf fibers through every
step.  A physical permutation changes the sides reading a node's flags;
the final histogram still counts each leaf in its original canonical
coordinates.  To apply Lemma~\ref{alg:orientation-banks}, let \(\tau_r\)
be the relative physical permutation of original leaf \(r\).  The exact
full-word law \(P\) fixes the count of canonical label \(i\) in relative
orientation \(\tau\) to be
\[
 H_{\tau,i}
 =N\sum_w P(w)\,
       \#\{r:\tau_r=\tau,\ w_r=i\}.
\]
These counts are independent of the direct output.  The lemma therefore
supplies equal original orientation banks when this schedule is taken in
all six global orientations.  No individual occurrence is required to carry
a fixed letter.
\end{proof}

To spell out the three-position case used in the library, let \(c(i,j)\)
be the two-child completion label: it is zero for \((i,j)=(0,0)\), and
is \(a\in\{1,2\}\) if both entries belong to \(\{0,a\}\) and at
least one equals \(a\).  It is undefined on \((1,2),(2,1)\).  Write
\[
 a=c(i,j),\qquad z=a-k\pmod3,\qquad s=c(z,r),
 \quad k\in\{0,1,2\}.
\]
The subtraction by \(k\) represents the physical rotation of the whole
first completed block.  In a fixed final fiber \(s\), first resolve the
top-level operand labels \((z,r)\).  This fixes \(a=z+k\pmod3\),
after which the inner word \((i,j)\) can be resolved using the pair of
sides belonging to that inner completion, with its physical rotation
included.  If \(C_s\) is an arbitrary rational law on the retained
triples, the two levels supply
\begin{equation}
 H(C_s)=H(Z,R\mid S=s)
       +\sum_{z,r}\Prob(Z=z,R=r\mid S=s)
                      H(I,J\mid Z=z,R=r,S=s).
 \label{alg:nested-completion-entropy}
\end{equation}
Here \(Z\) is a deterministic function of \(I,J\), so no intermediate
label adds a second entropy count.  The intermediate rotation changes
the grouping flag and its physical reader; it does not recode \(I,J,R\)
when the original canonical word histogram is calculated.

\begin{proposition}[Composition of finite recipe laws]
\label{alg:recipe-composition}
Consider a primitive completion, a finite tree of completion gates as in
Lemma~\ref{alg:completion-hierarchy}, or a parity operation on \(b\)
original occurrences.
Let \(s\) be its outer W label, with a rational law \(a_s\), separated
by a fixed recipe tree with limiting score \(u_0\).  In the outer group \(s\),
let \(t\) be a further W label, with conditional rational law \(r_{st}\),
separated by a fixed recipe tree with limiting score \(u_s\).
For completion this label is deterministic and \(u_s=0\).
For parity it is the other of the majority and odd-count labels.
For each permitted \((s,t)\), choose a rational distribution \(P_{st}\)
on its remaining original coarse words.  If \(n_i(w)\) counts canonical
input label \(i\) in the word \(w\), then the resulting original law and
limiting score are
\begin{align}
 q_i&=\frac1b\sum_{s,t}a_s r_{st}\,
                   \mathbb E_{w\sim P_{st}}n_i(w),
       \label{alg:recipe-law}\\
 u&=\frac1b\left(u_0+\sum_s a_su_s
                     +\sum_{s,t}a_s r_{st}H(P_{st})\right).
       \label{alg:recipe-score}
\end{align}
The formula is feasible for every stipulated finite collection of child
recipes and rational distributions; no optimality assertion is required.
\end{proposition}
\begin{proof}
Use the outer child recipe while the internal choices remain unresolved
fibers.  Its coordinate-preserving identity is valid by
Lemma~\ref{alg:payload}.  Once its label is direct, apply the corresponding
conditional child in each group.  The group populations are fixed by exact
types, so the prescribed auxiliary product is appended once and shared by
all branches.  After these labels are direct, resolve the remaining words
by the hierarchy in Lemma~\ref{alg:completion-hierarchy} for completion,
or by the single remaining pairwise-readable placement in
Lemma~\ref{alg:parity}.  Prescribe the conditional types induced by
\(P_{st}\) at every refinement.  Their conditional entropies sum to
\(H(P_{st})\), and Lemma~\ref{lem:group} supplies exactly this limiting
score across the hierarchy.  The relative population of that pool is
\(a_sr_{st}\).

The same types provide the histograms required for banking.  The outer and
conditional child recipes fix the number of macro occurrences assigned to
each \((s,t)\) pool in each of their original banks.  Within every such
pool, impose the full-word law \(P_{st}\) through its readable refinements.
For each original leaf position, this law fixes its marginal counts across
the pool, just as in the preceding proof.  Grouping those contributions by
the leaf's relative physical permutation fixes the total histograms
\(H_\tau\).  This uses each child's whole-bank histogram and the new
full-word types, not a fixed letter at an individual position of a child
recipe.  Lemma~\ref{alg:orientation-banks} now makes the final canonical
law common to all six original banks.

Logarithms of counts, auxiliary rank bounds, and pairing volumes add.
Dividing their sum by the number \(b\) of original occurrences per new
block proves \eqref{alg:recipe-score}.  Counting the original labels in
the same conditional pools proves \eqref{alg:recipe-law}.  A finite common
multiple of the rational denominators and the fixed child bank sizes
makes all pool sizes integral.  Increasing these supplies makes the finitely
many group and type errors vanish.  Every counted original coarse word
survives the stipulated support restrictions, so the outputs on the scalar
W input are nonzero.  Arbitrary unresolved fibers are preserved
coefficientwise; a zero attached fiber remains zero.
\end{proof}

The words in \eqref{alg:recipe-law} are counted in their original canonical
coordinates.  The resulting \(q\) is the exact common law of the original
orientation banks, providing the finite input to the banked-value
construction of the next section.
